%------------------------------------------------------------------------------
\subsection{\file{models/kazmierczak2024/model.jl}: the K24 model, its options and its sliding laws}\label{ref:k24model}
\emph{Tour PART 12.}

The largest declarative file: every option type of Section~\ref{arch:options}, the seven sliding laws, the routing tape,
the parameter and workspace structs, the model and its constructor. Its docstrings carry most of the design history
(why each default is what it is); the essentials are below.

\subsubsection*{Option types}
\begin{fnlist}
\item[AbstractDissipationMelt, DissipationMeltOn, DissipationMeltOff] Whether $Q_{\rm diss}$ enters the source. The paper drops it as
  negligible; the code keeps it on by default, ``for generality''. With it off and an $N$-independent law, the solve is a single
  routing pass.
\item[AbstractPsiOutAlgorithm, RecursivePsiOut, IterativePsiOut, TopologicalPsiOut(allow\_cycles), TapedPsiOut] How the routing graph is
  traversed (Section~\ref{sec:routing}). Recursive and iterative give the same field bit for bit; taped replays a recorded traversal
  and also matches bit for bit; topological is exact only on an acyclic graph and errors on a cycle unless \code{allow\_cycles = true}.
  The keyword constructor \code{TopologicalPsiOut(; allow\_cycles = false)}. Default: \code{TapedPsiOut()}; the other three implement
  only the GDSWarner routing.
\item[AbstractFrictionDiscretization, CellCentredFriction, StaggeredFriction\{M\}(ux, uy, quadrature, u\_floor)] Section~\ref{disc:friction}. The
  outer constructor \code{StaggeredFriction(ux, uy; quadrature = false, u\_floor = perYear2perSecond(1e-3))} stores the velocities as
  \code{Matrix\{Float64\}} copies; a coupled driver overwrites them in place each step.
\item[AbstractRoutingScheme, GDSWarner, Warner, Quinn(original), Tarboton, ModifiedTarboton, GDSTarboton] Section~\ref{disc:schemes}.
  \code{Quinn(; original = false)}.
\item[n\_directions(scheme)] 4 for \code{GDSWarner} and \code{Warner}, 8 otherwise. Face-based options need 4.
\item[AbstractQConversion, QFromOutflow, QFromFaceAverage] Section~\ref{disc:qconv}.
\item[AbstractDissipationDiscretization, CellCentredDissipation, FaceDissipation] Section~\ref{disc:dissipation}.
\item[AbstractFillAlgorithm, JacobiFill, LowestNeighbourFill, PriorityFloodFill(epsilon)] Section~\ref{disc:filling};
  \code{PriorityFloodFill(; epsilon = 1.0)} [Pa].
\item[RoutingTape\{T\}] Mutable. The recorded traversal: operation $k$ is \code{psi\_out[dst] += psi\_out[src] * w[k]}, or the clamp
  \code{psi\_out[dst] = max(0, psi\_out[dst])} when \code{src\_i[k] == 0} (indices stored as \code{Int32}); \code{valid}; and the arrays shared
  by the non-GDSWarner schemes: \code{w8} $(8, N_x, N_y)$ outflow fractions, face fluxes \code{Fx} $(N_x{+}1, N_y)$ and \code{Fy}
  $(N_x, N_y{+}1)$ [m$^3$/s], and \code{diss}, the face-assembled dissipation melt [kg\,m$^{-2}$\,s$^{-1}$].
  \code{RoutingTape\{T\}(Nx, Ny)} builds an empty, invalid tape.
\item[AbstractWaterThicknessAlgorithm, ArealConduitThickness, DarcyWeisbachThickness\{G\}, LaminarThickness\{G\}] Section~\ref{phys:W}. The
  docstring explains why $H$ (the depth inside one conduit, not an areal quantity) is deliberately not offered as \code{state.W}. Keyword
  constructors \code{DarcyWeisbachThickness(; gradient\_convention = MeanGradient())}, likewise \code{LaminarThickness}.
\item[AbstractGradientConvention, LocalGradient, MeanGradient] Local or grounded-mean gradient in the two sheet closures.
\item[AbstractDrainageMode, BothDrainage, EfficientOnly, InefficientOnly] Section~\ref{phys:drainage}.
\end{fnlist}

\subsubsection*{Sliding laws}
\begin{fnlist}
\item[AbstractSlidingLaw, AbstractPressureDependentSlidingLaw] The second marks laws whose $\tau_b$ depends on $N$, which makes
  \code{resolve\_q!} take the joint $(q, N)$ loop.
\item[NoFrictionSlidingLaw] $\tau_b = 0$. The default: frictional heat is opt-in.
\item[PrescribedFieldSlidingLaw\{A\}(tau\_b), PrescribedFieldSlidingLaw(grid, tau\_b)] A fixed scalar $|\tau_b|$ field (e.g.\ Yelmo's
  \code{taub}, or $\sqrt{\tau_x^2+\tau_y^2}$), wrapped with \code{alloc\_field}. Mirrors Shakti's \code{PrescribedSlidingLaw}; the docstring
  argues a real $\tau_b$ field is better grounded than a guessed Coulomb coefficient, and it needs no $(q,N)$ loop. This is what the
  coupled Yelmo runs use.
\item[WeertmanSlidingLaw\{T\}(C, q), WeertmanSlidingLaw(; C, q = 1/3)] $C|u_b|^q$.
\item[PowerPlasticSlidingLaw\{T\}(c\_till, q, u0); keyword constructor] $c N (|u_b|/u_0)^q$; $q = 1$, $u_0 = 100$\,m/yr; Yelmo's
  \code{\_calc\_beta\_aa\_power\_plastic!}.
\item[RegularizedCoulombSlidingLaw\{T\}(c\_till, q, u0); keyword constructor] $c N (|u_b|/(|u_b|+u_0))^q$; $q = 1/3$; Yelmo's
  \code{\_calc\_beta\_aa\_reg\_coulomb!}.
\item[RegularizedCoulombFieldSlidingLaw\{A, F\}; (grid, c\_till; q, u0)] Same, with a per-cell $c$ (e.g.\ Yelmo's \code{cb\_ref}).
\item[ShaktiRegularizedCoulombSlidingLaw\{A, F\}(C, n, inv\_n, lambda); (grid, A\_visc, C; n = 3, lambda\_coeff = 1.5)] Shakti's own
  regularisation with $\lambda = 1.5A$ per cell, for bit-for-bit parity with a Shakti run. Has no \code{calc\_tau\_b} method, since
  $\lambda$ is a field.
\item[convert\_sliding\_law(T, law)] Converts a law's scalars to the model's float type; field laws are returned unchanged
  (\code{===}) when their element type already matches, converted otherwise.
\end{fnlist}

\subsubsection*{The model}
\begin{fnlist}
\item[KazmierczakParams\{T, D, L, P, WT, M\}] 39 immutable fields: the physical constants (\code{rho\_w} \ldots \code{eta\_w}), the clamps
  (\code{Wmin}, \code{Wmax}, \code{sigmat}, \code{q\_min}, \code{q\_max}), \code{longcoupwater}, every option, and the iteration controls
  (\code{fill\_iters}, \code{max\_psi\_out\_calls}, \code{max\_dissipation\_iters}, \code{dissipation\_rtol}, \code{dissipation\_verbose},
  \code{max\_qN\_iters}, \code{qN\_rtol}, \code{qN\_verbose}). Type parameters: float type, dissipation trait, sliding law, psi-out algorithm,
  thickness closure, drainage mode.
\item[KazmierczakWorkspace\{A, R\}] 38 fields, all allocated once: the potential and its gradients (9 fields), the routing and source
  fields (\code{visited}, \code{h}, \code{G}, \code{q\_T}, \code{i\_eb}, \code{Q\_b}, \code{Q\_diss}, \code{mdot\_fixed}, \code{mdot\_total}, \code{psi\_out},
  \code{corfac}, \code{q}, \code{q\_prev}, \code{tau\_b}, \code{N\_prev}), the \code{routing\_tape} and three filling scratch buffers, and the
  conduit fields (\code{Q}, \code{kappa}, \code{abs\_v\_b}, \code{A\_visc}, \code{S\_inf}, \code{H\_hard}, \code{H\_soft}, \code{H}, \code{N\_inf},
  \code{Po}). Appendix~\ref{app:fields} lists them all.
\item[KazmierczakHydroModel\{T, A, D, L, P, WT, M\}(params, workspace)] Section~\ref{arch:split}.
\item[getproperty(model, name), propertynames(model)] Forward \code{model.<name>} to \code{params}, then \code{workspace}.
\item[KAZMIERCZAK\_DEFAULT\_L\_W] $3.34\times10^5$ J/kg; the tests use it to supply a melt rate as $G = \mdot L_w$.
\item[KazmierczakHydroModel(grid, kappa, abs\_v\_b, A\_visc, G, q\_T; kwargs...)] The constructor. In order: checks every input (and
  \code{i\_eb}) is \code{(Nx, Ny)}; converts every scalar to \code{T = typeof(grid.dx)}; computes $K$; handles
  \code{coupling\_length\_kamb86} (unset: a long \code{@warn} explaining how to choose it from the grid resolution, then 10; negative:
  \code{ArgumentError}; \code{longcoupwater} = value/2); fills the unset \code{fill\_algorithm}, \code{q\_conversion} and
  \code{dissipation\_discretization} from the routing scheme; rejects a non-GDSWarner scheme without \code{TapedPsiOut}, a
  \code{StaggeredFriction} of the wrong size, and face-based options with an 8-neighbour scheme; turns \code{dissipation\_melt} into the
  trait; converts the sliding law; allocates the workspace and sets \code{mdot\_fixed} $= (G - q_T)/L_w$. All keywords are listed in
  Section~\ref{opt:k24}.
\item[set\_basal\_terms!(model; G, q\_T, i\_eb)] Replaces any of the given terms in place (a \code{nothing} keyword keeps the
  current value) and recomputes \code{mdot\_fixed}. The only way to change the source: the model never takes a melt rate, so a
  coupled model exchanges heat fluxes and water, not melt.
\end{fnlist}

%------------------------------------------------------------------------------
\subsection{\file{models/kazmierczak2024/sliding\_law.jl}: $\tau_b$ and the frictional heat}\label{ref:sliding}
\emph{Tour PART 13.}
\begin{fnlist}
\item[calc\_tau\_b(law, N, abs\_v\_b)] The scalar formula of each $N$-free or scalar-coefficient law (NoFriction, Weertman,
  PowerPlastic, RegularizedCoulomb), for tests and diagnostics.
\item[update\_tau\_b!(model, state, law)] Writes \code{model.tau\_b} from \code{state.N} and \code{model.abs\_v\_b}, one method per law. The
  formulas are written out again, with the law's parameters pulled into local scalars, rather than calling \code{calc\_tau\_b} inside the
  broadcast, so the broadcast is plain arithmetic on fields and scalars. RegularizedCoulomb and its field variant share one method
  (broadcasting handles a scalar or field $c$); the tour checks they agree exactly for a uniform field.
\item[add\_friction\_term!(model, fd::StaggeredFriction)] $Q_b$ from the C-grid velocities (\code{staggered\_friction\_kernel!}), then
  \code{mdot\_total += Q\_b/L\_w}. (The generic and cell-centred methods are in \file{water\_flux.jl}.)
\item[stagger\_beta(tau\_b, abs\_v\_b, u\_floor, i, j)] $\beta = \tau_b/\max(|u_b|, u_{\rm floor})$ at a centre.
\item[GQ\_S3, GQ\_PTS, gq\_interp(c, p)] The $2\times2$ Gauss points $(\pm1/\sqrt3, \pm1/\sqrt3)$ of the reference cell and bilinear
  interpolation from the SW, SE, NE, NW corner values (Yelmo's \code{gq2D}).
\item[acx\_corners(F, i, j, im1, jm1, jp1), acy\_corners(F, i, j, im1, ip1, jm1)] Corner (ab-node) values of an acx/acy field around cell
  $(i,j)$, each the mean of the two faces meeting there.
\item[staggered\_friction\_kernel!(Q\_b, tau\_b, abs\_v\_b, ux, uy, Nx, Ny, u\_floor, quadrature)] First the face tractions
  $\tau_x = \tfrac12(\beta_i + \beta_{i+1})u_x$, $\tau_y$ likewise (edge faces clamp the neighbour index), then either the face
  products averaged to the centre or the Gauss-point average of $|\mathbf u||\boldsymbol\tau|$. Allocates two scratch arrays per call.
\end{fnlist}
